\section{Conclusion}\label{sec:conclusion}

This study transforms a GA-based conference precursor into a controlled Neural Computing and Applications benchmark for neural intraday trend classification. It preserves the predecessor's practical insight---a compact, Information-Gain-selected technical representation can support metaheuristic hyperparameter search---but replaces randomized validation and a single-optimizer design with a chronological dual-fitness protocol, five budget-matched optimizers, and locked predictive and economic assessment.

The current evidence indicates that the choice between compound MCC/$F_1$ fitness and accuracy fitness produces configuration-dependent changes in both held-out accuracy and MCC. Within the accuracy-fitness protocol, MLP is the strongest model in the exploratory predictive analysis and leads mean total profit, although RF has the smallest mean drawdown. These results reinforce the central methodological lesson: selection objective, reported metric, and economic utility must be treated as related but distinct quantities.

The manuscript deliberately avoids a confirmatory claim about a universal best fitness function because only three paired seeds are currently available across every cell. Completing the common 20-seed design is the necessary next step. The resulting framework nevertheless provides a transparent and reproducible basis for future optimizer comparisons in non-stationary financial time series.
